\chapter{Introduzione}


\section{Descrizione}
Il progetto \textbf{ChargeNet} è un sistema software pensato per la gestione di una rete di stazioni di ricarica per veicoli elettrici (EV). L'idea alla base dell'applicazione è quella di coordinare tutto ciò che ruota attorno a una sessione di ricarica: dall'aiutare l'utente a trovare una colonnina libera, fino alla gestione del pagamento, assicurandosi allo stesso tempo che l'infrastruttura elettrica non subisca danni.

Per raggiungere questo obiettivo, la piattaforma fa da punto di incontro per tre specifiche figure, fornendo a ciascuna le funzionalità necessarie:

\begin{itemize}
    \item \textbf{Driver:} È l'utente finale. Utilizza il sistema per cercare le stazioni compatibili con la sua auto, far partire e monitorare la ricarica, e gestire i pagamenti tramite un portafoglio virtuale interno (wallet).
    
    \item \textbf{Station Operator:} È l'addetto alla manutenzione fisica. Il suo compito è controllare lo stato di salute delle colonnine: se una stazione si guasta, può metterla temporaneamente fuori servizio per evitare disagi ai Driver, e registrarla di nuovo come attiva una volta riparata.
    
    \item \textbf{Energy Manager:} È il tecnico supervisore della rete elettrica. Controlla in tempo reale il carico sui trasformatori per evitare surriscaldamenti. Se la rete è sotto stress, ha gli strumenti per abbassare la potenza erogata o limitare le ricariche (Load Balancing) in modo da proteggere l'infrastruttura.
\end{itemize}

\section{Stack Tecnologico}
Per lo sviluppo è stato scelto un ecosistema Java-centrico, privilegiando 
la robustezza e la portabilità:

\begin{itemize}
    \item \textbf{Linguaggio:} Java 21.
    \item \textbf{Interfaccia Grafica:} JavaFX 21, per una gestione 
    reattiva delle dashboard dei vari attori.
    \item \textbf{Persistenza:} PostgreSQL, gestito tramite JDBC per un 
    controllo granulare delle transazioni.
    \item \textbf{Build Tool:} Maven, per la gestione delle dipendenze e 
    del ciclo di vita del software.
\end{itemize}

\section{Astrazione dell'Infrastruttura Fisica}
Trattandosi di un progetto universitario, l'applicazione non comunica con colonnine e trasformatori reali. Nel mondo reale, infatti, interfacciarsi con l'hardware fisico è un'operazione che porta con sé diverse complicazioni: bisogna gestire le differenze tra i vari produttori, adattarsi a protocolli di rete specifici e gestire i fisiologici ritardi o cadute di connessione dei sensori.

L'architettura del sistema è stata comunque pensata per un'eventuale integrazione futura. Per le finalità di questo progetto, i dati sui consumi e sullo stato delle stazioni vengono semplicemente simulati da un componente interno al programma (\textit{GridMonitor}). 